\begin{resumo}

Este trabalho estima ganhos de proficiência associados à adesão ao Programa Ensino Integral (PEI), política que ampliou a jornada escolar e reorganizou a gestão das escolas estaduais em São Paulo. A pesquisa explora a adoção escalonada do programa a partir de um arcabouço de diferença-em-diferenças, comparando Efeitos Fixos de Duas Vias (TWFE), \textcite{callawaysantanna}, \textcite{santannazhao} e uma aproximação agregada inspirada em \textcite{santannaxu}. As estimativas principais indicam ganhos médios pós-tratamento entre aproximadamente 0,50 e 0,68 desvio-padrão em exames padronizados de Língua Portuguesa e Matemática. A incorporação de covariáveis de composição construídas a partir dos microdados de matrícula e de controles municipais não altera substancialmente as estimativas, indicando que a composição observável agregada não é suficiente, nas especificações estimadas, para eliminar os ganhos. A interpretação, contudo, é limitada: a ausência de vínculo individual entre matrícula e nota impede decompor integralmente aprendizado, seleção e participação na prova, de modo que os resultados devem ser lidos como evidência agregada condicionada às hipóteses de identificação adotadas.

\textbf{Palavras-chave:} \imprimirpalavraschavesresumo.
\end{resumo}
